\section{Síntesis y ampliación}

Lo que esta clase muestra en realidad es un data-to-model control loop: primero se usan el scaling y el lifecycle cost para decidir la escala objetivo; luego se eligen las fuentes y la provenance; después se buscan los filtros con ablation, se atienden opt-out/PII/license con governance, se define la tarea con formatting/mixture y, por último, se revisa el modelo con contamination-aware evaluation. El valor de StarCoder no está solo en el code score: está en que toda esta cadena puede auditarse.

\subsection{Cadena causal de extremo a extremo}

\begin{importantbox}{De la data decision al model behavior}
\begin{enumerate}
  \item \textbf{Source} determina el conocimiento, las licencias y la cobertura de la cola larga;
  \item \textbf{Filter/dedup} determina la calidad, la diversity y la memorization;
  \item \textbf{Formatting/mixture} determina las tareas y el contexto que ve el modelo;
  \item \textbf{Training scale} determina la capacidad y la exposure;
  \item \textbf{Evaluation} determina si podemos distinguir la generalización de la contamination.
\end{enumerate}
Si cualquiera de estas capas no es visible, el benchmark final es difícil de interpretar.
\end{importantbox}

\subsection{Los cinco errores más comunes}

\begin{warningbox}{Del headline a una conclusión falsa}
No presentes la Chinchilla ratio como una constante válida entre domains; no uses stars/comments como quality label; no presentes un PII detector como prueba de riesgo cero; no presentes la puntuación de HumanEval como capacidad de ingeniería de software; no presentes los open weights como openness completa. En el curso, los resultados negativos y las limitaciones de gobernanza importan tanto como los resultados positivos.
\end{warningbox}

\subsection{Lista de verificación para reproducir los experimentos}

\begin{lstlisting}[caption={Lista unificada de reproducción para proyectos de preentrenamiento}]
pin_model_tokenizer_optimizer_and_code_commit()
hash_raw_filtered_and_formatted_dataset_versions()
report_source_mixture_repeats_and_effective_tokens()
publish_filter_ablations_seeds_and_negative_results()
document_opt_out_pii_license_and_decontamination()
freeze_evaluation_prompts_harnesses_cutoffs_and_dates()
report_train_and_inference_compute_memory_latency_cost()
\end{lstlisting}

\subsection{Preguntas de autoevaluación}

\begin{multicols}{2}
\footnotesize
\begin{enumerate}
  \setlength{\itemsep}{0pt}
  \item ¿Por qué los open weights no equivalen a un transparent release completo?
  \item ¿En qué difieren Kaplan y Chinchilla respecto de la model/data allocation?
  \item ¿Por qué lo compute-optimal no es necesariamente lifecycle-optimal?
  \item ¿Qué diferencia hay entre raw bytes, unique tokens y training tokens?
  \item ¿Por qué no se puede tokenizar Common Crawl directamente y entrenar con él?
  \item ¿Qué data lineage requiere el opt-out de The Stack v1?
  \item ¿Qué cambia en v2 el uso de Software Heritage?
  \item En los synthetic data, ¿qué controla cada uno: el seed, el generator y el verifier?
  \item ¿Cómo aumenta Cosmopedia la diversidad de lo generado?
  \item ¿Por qué el repository stars filter produjo la peor ablation?
  \item ¿Para qué sirven varios seeds en un filter experiment?
  \item ¿Por qué MinHash/LSH es adecuado para la near-dedup?
  \item ¿Por qué una strong dedup puede mejorar un code benchmark?
  \item ¿Qué diferencia hay entre PII detection y secret scanning?
  \item ¿Por qué la decontamination debe abarcar los SFT data?
  \item ¿Cómo cambian los metadatos de repository/file la tarea de entrenamiento?
  \item ¿Cómo se implementa FIM con un next-token objective?
  \item ¿Cómo cambia el source sampling weight la effective mixture?
  \item ¿Qué optimizan principalmente MQA y GQA?
  \item ¿Por qué la openness de BigCode va más allá de publicar los pesos?
  \item ¿Pueden compararse directamente pass@1 y pass@k?
  \item ¿Por qué un membership test no puede demostrar memorization?
  \item ¿Qué tipo de fuga reduce la time split de LiveCodeBench?
  \item Para construir un code assistant con una sola GPU, ¿en qué conviene invertir primero?
\end{enumerate}
\end{multicols}

\subsection{Lecturas adicionales}

\begin{multicols}{2}
\footnotesize
\begin{itemize}
  \setlength{\itemsep}{0.15em}
  \item Hoffmann et al., \href{https://arxiv.org/abs/2203.15556}{Training Compute-Optimal Large Language Models}: Chinchilla scaling.
  \item Muennighoff et al., \href{https://arxiv.org/abs/2305.16264}{Scaling Data-Constrained Language Models}: entrenamiento con repetición cuando los unique data son limitados.
  \item Kocetkov et al., \href{https://arxiv.org/abs/2211.15533}{The Stack}; Lozhkov et al., \href{https://arxiv.org/abs/2402.19173}{StarCoder 2 and The Stack v2}.
  \item Penedo et al., \href{https://arxiv.org/abs/2306.01116}{RefinedWeb} y \href{https://arxiv.org/abs/2406.17557}{FineWeb}: web filtering y ablations.
  \item Lee et al., \href{https://arxiv.org/abs/2107.06499}{Deduplicating Training Data Makes Language Models Better}.
  \item Li et al., \href{https://arxiv.org/abs/2305.06161}{StarCoder: may the source be with you!}: modelo, FIM, gobernanza y evaluación.
  \item Brown et al., \href{https://arxiv.org/abs/2107.03374}{HumanEval}; Jain et al., \href{https://arxiv.org/abs/2403.07974}{LiveCodeBench}: evaluación por execution y con separación temporal.
  \item Zhou et al., \href{https://arxiv.org/abs/2305.11206}{LIMA}; Hu et al., \href{https://arxiv.org/abs/2106.09685}{LoRA}: alignment con pocos datos de alta calidad y parameter-efficient fine-tuning.
\end{itemize}
\end{multicols}

\end{document}
